\documentclass[12pt]{article}

\usepackage{amsmath,amssymb,amsthm}
\usepackage{mathtools}
\usepackage{geometry}
\usepackage{setspace}
\usepackage{hyperref}
\usepackage{booktabs}
\usepackage{xcolor}

\geometry{letterpaper, margin=1in}
\setstretch{1.15}

% Theorem environments
\newtheorem{theorem}{Theorem}[section]
\newtheorem{lemma}[theorem]{Lemma}
\newtheorem{proposition}[theorem]{Proposition}
\newtheorem{corollary}[theorem]{Corollary}
\newtheorem{definition}[theorem]{Definition}
\newtheorem{remark}[theorem]{Remark}

% Shortcuts
\newcommand{\Z}{\mathbb{Z}}
\newcommand{\R}{\mathbb{R}}
\newcommand{\N}{\mathbb{N}}
\DeclareMathOperator{\csch}{csch}
\DeclareMathOperator{\sech}{sech}

\title{The Discrete--Continuous Bridge Through the Lemniscatic Lens:\\[6pt]
\large An Analysis of $\displaystyle\sum_{n=0}^{\infty}\frac{1}{(n^2+1)^2}$}

\author{W.\,J.\ Steinmetz~III}
\date{}

\begin{document}
\maketitle

\begin{abstract}
We present a complete derivation of the closed-form identity
\[
\sum_{n=0}^{\infty}\frac{1}{(n^2+1)^2}
\;=\;\frac{2+\pi\coth\pi+\pi^2\csch^2\pi}{4}\,,
\]
then recast the result entirely in terms of the lemniscate constant~$\varpi$
and the bridge constant ${G^{*}}=2\varpi/\sqrt{\pi}$ of Foundational Ternary
Dynamics~(FTD).  In the FTD rewriting, the ontic content of the identity
reduces to the single ratio $\ell=\varpi^2/{G^{*}}^2=\pi/4$, revealing
that every appearance of~$\pi$ in the standard formula is a derived
consequence of this deeper lemniscatic ratio.

We show that the sum is a restriction of the Epstein zeta function
$Z_{\Z[i]}(s)$ at $s=2$, establish the exact identity
$E_4(i)=3\,{G^{*}}^8/(256\,\varpi^4)$, and prove the Epstein factorisation
$Z_{\Z[i]}(2)=4\,\zeta(2)\,L(2,\chi_{-4})=32\,\varpi^4\,G/(3\,{G^{*}}^4)$
in ontic variables, where $G$ is Catalan's constant---the irreducible
arithmetic residue encoding the distribution of Gaussian primes.

We identify the exponential scale $e^{4\ell}=e^{\pi}$ (Gelfond's
constant) as the natural decay parameter controlling the lattice's
discrete-to-continuous transition, and show that the row restriction
from~$\Z[i]$ to a single horizontal line eliminates the arithmetic factor,
explaining why our one-dimensional sum admits a purely ontic closed form.

We then place these results in a dimensional hierarchy: the Basel sum
$\zeta(2) = \pi^2/6$ is the lattice zeta function of~$\Z$ at $s = 2$;
the Eisenstein series $G_4$ encodes the period of~$\Z[i]$ through~$\varpi$;
the Watson integral encodes the period of~$\Z^3$ through~${G^{*}}$.
Each dimension has its own ``circle constant,'' and only $D = 3$
produces a self-consistent framework.

This paper is one of several papers presenting Foundational Ternary
Dynamics.
\end{abstract}

\tableofcontents

%=======================================================================
\section{The Problem}
%=======================================================================

We seek a closed form for the sum
\begin{equation}\label{eq:S}
S \;=\; \sum_{n=0}^{\infty}\frac{1}{(n^2+1)^2}\,.
\end{equation}
This is a convergent series of rational terms indexed by the non-negative
integers.  Its evaluation requires passing from the discrete lattice~$\N$
into the analytic structure of the complex plane, and the result is a
transcendental number expressed in terms of hyperbolic functions of~$\pi$.
The purpose of this note is to make the architecture of that passage
explicit and to identify its ontic content.

%=======================================================================
\section{Standard Derivation}
%=======================================================================

\subsection{Strategy}

The exponent~2 in the denominator of~\eqref{eq:S} makes direct summation
difficult.  We proceed by first evaluating the simpler sum
$\sum 1/(n^2+a)$ in closed form, then recovering~$S$ by differentiation
with respect to the auxiliary parameter~$a$.

\subsection{The Bilateral Sum}

The Mittag-Leffler expansion of $\pi\cot(\pi z)$ gives the classical identity
\begin{equation}\label{eq:bilateral}
\sum_{n=-\infty}^{\infty}\frac{1}{n^2+a}
\;=\;\frac{\pi}{\sqrt{a}}\,\coth(\pi\sqrt{a})\,,
\qquad a>0.
\end{equation}
This can be derived from the partial fraction decomposition of
$\pi\cot(\pi z)$, or equivalently from the Poisson summation formula
applied to the Lorentzian $f(x)=1/(x^2+a)$.

\begin{remark}[First bridge: lattice to analysis]\label{rem:bridge1}
Equation~\eqref{eq:bilateral} is already a discrete-to-continuous bridge:
the left side is a sum over the integer lattice~$\Z$, while the right side
is a smooth function of the continuous parameter~$a$, expressed in terms of
the hyperbolic cotangent.  The Poisson summation route makes the mechanism
explicit: the Fourier transform converts the sum over direct-lattice points
into a sum over reciprocal-lattice points, and the latter is dominated by a
single exponential scale.
\end{remark}

\subsection{The Unilateral Sum}

Since $1/(n^2+a)$ is symmetric in~$n$, the bilateral sum double-counts
every term except the $n=0$ term.  Hence
\begin{equation}\label{eq:unilateral}
\sum_{n=0}^{\infty}\frac{1}{n^2+a}
\;=\;\frac{1}{2}\!\left(\frac{\pi}{\sqrt{a}}\,\coth(\pi\sqrt{a})+\frac{1}{a}\right)
\;=\;\frac{\pi\coth(\pi\sqrt{a})}{2\sqrt{a}}+\frac{1}{2a}\,.
\end{equation}
Note that every term on the left is positive, and the right side is
manifestly positive for $a>0$, as required.

\subsection{Differentiation with Respect to the Parameter}

Observe that
\begin{equation}\label{eq:diff-trick}
\frac{1}{(n^2+a)^2}\;=\;-\frac{\partial}{\partial a}\!\left(\frac{1}{n^2+a}\right).
\end{equation}
Summing over $n\ge 0$ and exchanging the sum and derivative (justified by
uniform convergence on compact subsets of $(0,\infty)$):
\begin{equation}\label{eq:sum-diff}
\sum_{n=0}^{\infty}\frac{1}{(n^2+a)^2}\;=\;-\frac{d}{da}\sum_{n=0}^{\infty}\frac{1}{n^2+a}\,.
\end{equation}
Define
\begin{equation}\label{eq:g}
g(a)\;=\;\frac{\pi\coth(\pi\sqrt{a})}{2\sqrt{a}}+\frac{1}{2a}\,.
\end{equation}
Then
\begin{equation}\label{eq:S-gprime}
S\;=\;-\,g'(1).
\end{equation}

\begin{remark}[Second bridge: discrete to continuous]\label{rem:bridge2}
The fixed problem at $a=1$ has been embedded into a one-parameter
continuous family.  The answer is recovered by differentiation followed by
evaluation, converting a discrete summation problem into a calculus problem.
\end{remark}

\subsection{Computing $g'(a)$}

Write $g=g_1+g_2$ with
\[
g_1(a)=\frac{\pi\coth(\pi\sqrt{a})}{2\sqrt{a}}\,,\qquad
g_2(a)=\frac{1}{2a}\,.
\]
For $g_2$:
\begin{equation}\label{eq:g2prime}
g_2'(a)=-\frac{1}{2a^2}\,.
\end{equation}
For $g_1$, apply the product rule to
$\frac{\pi}{2}\cdot\coth(\pi\sqrt{a})\cdot a^{-1/2}$:
\begin{align}
g_1'(a)
&=\frac{\pi}{2}\left[\frac{d}{da}\bigl(\coth(\pi\sqrt{a})\bigr)\cdot a^{-1/2}
  +\coth(\pi\sqrt{a})\cdot\bigl(-\tfrac{1}{2}\bigr)a^{-3/2}\right]\notag\\
&=\frac{\pi}{2}\left[\bigl(-\csch^2(\pi\sqrt{a})\bigr)\cdot
  \frac{\pi}{2\sqrt{a}}\cdot\frac{1}{\sqrt{a}}
  -\frac{\coth(\pi\sqrt{a})}{2a^{3/2}}\right]\notag\\
&=-\frac{\pi^2\csch^2(\pi\sqrt{a})}{4a}
  -\frac{\pi\coth(\pi\sqrt{a})}{4a^{3/2}}\,.\label{eq:g1prime}
\end{align}
Combining:
\begin{equation}\label{eq:gprime}
g'(a)\;=\;
-\frac{\pi^2\csch^2(\pi\sqrt{a})}{4a}
-\frac{\pi\coth(\pi\sqrt{a})}{4a^{3/2}}
-\frac{1}{2a^2}\,.
\end{equation}

\subsection{Evaluation at $a=1$}

Setting $a=1$ in~\eqref{eq:gprime}:
\[
g'(1)=-\frac{\pi^2\csch^2\pi}{4}-\frac{\pi\coth\pi}{4}-\frac{1}{2}
=-\frac{\pi^2\csch^2\pi+\pi\coth\pi+2}{4}\,.
\]
By~\eqref{eq:S-gprime}:
\begin{equation}\label{eq:result}
\boxed{
S\;=\;\sum_{n=0}^{\infty}\frac{1}{(n^2+1)^2}
\;=\;\frac{2+\pi\coth\pi+\pi^2\csch^2\pi}{4}
\;\approx\;1.30683697542\,.
}
\end{equation}

%=======================================================================
\section{FTD Recasting: Eliminating $\pi$}\label{sec:ftd}
%=======================================================================

\subsection{The Ontic Constants}

\begin{definition}[Lemniscate constant]\label{def:varpi}
The lemniscate constant~$\varpi$ is defined by
\[
\varpi \;=\; 2\int_0^1\frac{dt}{\sqrt{1-t^4}}
\;=\;\frac{\Gamma(1/4)^2}{2\sqrt{2\pi}}
\;\approx\;2.62206\,.
\]
It is the half-perimeter of the lemniscate of Bernoulli and satisfies
$\varpi=\pi/M(1,\sqrt{2})$, where $M$ denotes the arithmetic--geometric mean.
\end{definition}

\begin{definition}[Bridge constant]\label{def:Gstar}
The FTD bridge constant is
\[
{G^{*}} \;=\; \frac{\Gamma(1/4)^2}{\sqrt{2}\;\Gamma(1/2)^2} \;=\; \frac{2\varpi}{\sqrt{\pi}} \;\approx\; 2.95868\,.
\]
The first form is $\pi$-free: both $\Gamma(1/4)$ and $\Gamma(1/2)$ are definable by $\pi$-free integrals ($\int_0^\infty t^{s-1}e^{-t}\,dt$ at $s = 1/4$ and $s = 1/2$). The equivalence of the two forms follows from $\Gamma(1/2)^2 = \pi$ (the Gaussian integral).
\end{definition}

\begin{definition}[Ontic ratio]\label{def:ell}
Define
\begin{equation}\label{eq:ell}
\ell \;=\; \frac{\varpi^2}{{G^{*}}^2}\,.
\end{equation}
\end{definition}

\begin{lemma}\label{lem:pi}
$\pi = 4\ell = 4\varpi^2/{G^{*}}^2$.
\end{lemma}
\begin{proof}
From ${G^{*}}=2\varpi/\sqrt{\pi}$ we obtain ${G^{*}}^2=4\varpi^2/\pi$,
hence $\pi=4\varpi^2/{G^{*}}^2=4\ell$.
\end{proof}

\begin{remark}\label{rem:ontic-priority}
In the FTD framework, $\varpi$ and ${G^{*}}$ (equivalently $\varpi$ and
$\ell$) are treated as the ontic (structurally primary) constants.
The number $\pi$ is a derived quantity: it is four times the ratio
$\varpi^2/{G^{*}}^2$.  This reflects the claim that the lemniscatic lattice
structure is ontologically prior to the circular one.
\end{remark}

\subsection{The Sum in Ontic Variables}

\begin{theorem}[FTD form of the sum]\label{thm:ftd}
\begin{equation}\label{eq:ftd-form}
\boxed{
\sum_{n=0}^{\infty}\frac{1}{(n^2+1)^2}
\;=\;\frac{2+4\ell\coth(4\ell)+16\ell^2\csch^2(4\ell)}{4}\,,
}
\end{equation}
where $\ell=\varpi^2/{G^{*}}^2$.
\end{theorem}
\begin{proof}
Substitute $\pi=4\ell$ into~\eqref{eq:result}.
\end{proof}

\begin{remark}\label{rem:pi-disguised}
Every appearance of $\pi$ in the standard result was a disguised occurrence
of $4\varpi^2/{G^{*}}^2$.  The hyperbolic functions $\coth$ and $\csch^2$
depend on a single transcendental argument: $4\ell=4\varpi^2/{G^{*}}^2$.
\end{remark}

\subsection{Gelfond's Constant as the Exponential Scale}

\begin{proposition}\label{prop:gelfond}
The exponential scale governing the hyperbolic functions
in~\eqref{eq:ftd-form} is
\[
e^{4\ell}\;=\;e^{4\varpi^2/{G^{*}}^2}\;=\;e^{\pi}\;\approx\;23.1407\,,
\]
which is Gelfond's constant, $(-1)^{-i}$.
\end{proposition}
\begin{proof}
$4\ell=\pi$ by Lemma~\ref{lem:pi}; $e^{\pi}=(-1)^{-i}$ by Euler's
identity.
\end{proof}

The nome of the Gaussian lattice $\Z[i]$ (with respect to the modular
parameter $\tau=i$) is
\[
q\;=\;e^{2\pi i\tau}\Big|_{\tau=i}\;=\;e^{-2\pi}\;=\;
e^{-8\varpi^2/{G^{*}}^2}\;\approx\;1.87\times 10^{-3}\,.
\]
The smallness of~$q$ (controlled entirely by $\varpi^2/{G^{*}}^2$) is
responsible for the rapid exponential suppression that makes the $\csch^2$
correction contribute only ${\sim}1.4\%$ of the total sum, as we quantify
in Section~\ref{sec:bridge-structure}.

%=======================================================================
\section{Connection to the Gaussian Lattice}\label{sec:gaussian}
%=======================================================================

\subsection{The Sum as a Lattice Restriction}

Observe that $n^2+1=|n+i|^2$ for $n\in\Z$.  Therefore:
\begin{equation}\label{eq:S-lattice}
S\;=\;\sum_{n=0}^{\infty}\frac{1}{|n+i|^4}\,.
\end{equation}
The natural two-dimensional object associated to the Gaussian lattice
$\Z[i]$ is the \textbf{Epstein zeta function}
\begin{equation}\label{eq:epstein}
Z_{\Z[i]}(s)\;=\;\sideset{}{'}\sum_{z\in\Z[i]}|z|^{-2s}
\;=\;\sideset{}{'}\sum_{(m,n)\in\Z^2}(m^2+n^2)^{-s}\,,
\end{equation}
where the prime denotes omission of $(m,n)=(0,0)$.  Our sum~$S$ is the
restriction of this lattice sum to the horizontal line $\mathrm{Im}(z)=1$
(that is, $m$ arbitrary with $n=1$), further restricted to $m\ge 0$.
Explicitly, the full $n=1$ row of $Z_{\Z[i]}(2)$ contributes
$\sum_{m=-\infty}^{\infty}(m^2+1)^{-2}=2S-1$, since the $m=0$ term equals~1
and the remaining terms pair symmetrically.

\begin{remark}\label{rem:epstein-vs-eisenstein}
Unlike the unnormalized Eisenstein series
$G_4(\tau)=\sum'(m+n\tau)^{-4}$, which involves the \emph{fourth power}
of a complex number, our sum involves the \emph{squared modulus} raised to
the second power: $|m+ni|^{-4}=(m^2+n^2)^{-2}$.  These are distinct
objects; $G_4$ is a holomorphic modular form, while the Epstein zeta
function is real-valued and not holomorphic.  The connection to modular
forms passes through the identity
$Z_{\Z[i]}(s)=4\,\zeta(s)\,L(s,\chi_{-4})$, where $\chi_{-4}$ is the
non-principal Dirichlet character modulo~4.
\end{remark}

\subsection{The Eisenstein Series in Ontic Variables}

Although our sum is a restriction of the Epstein zeta function rather than
of $G_4$ directly, the Eisenstein series at the Gaussian point still
controls the lattice's fundamental structure.  The normalized Eisenstein
series at $\tau=i$ is:
\[
E_4(i) \;=\; 1+240\sum_{n=1}^{\infty}\sigma_3(n)\,q^n\,,
\qquad q=e^{-2\pi}\,.
\]

\begin{theorem}[Eisenstein $E_4(i)$ in FTD form]\label{thm:E4}
\begin{equation}\label{eq:E4}
E_4(i)\;=\;3\!\left(\frac{\varpi}{\pi}\right)^{\!4}
\;=\;\frac{3\,{G^{*}}^8}{256\,\varpi^4}\,.
\end{equation}
\end{theorem}
\begin{proof}
The first equality is a classical result relating $E_4(i)$ to the
lemniscate constant via the Weierstrass invariant~$g_2$ of the square
lattice (see Zucker~\cite{zucker77}).  For the second equality, substitute
$\pi=4\varpi^2/{G^{*}}^2$:
\[
3\!\left(\frac{\varpi}{\pi}\right)^{\!4}
= 3\!\left(\frac{\varpi\,{G^{*}}^2}{4\varpi^2}\right)^{\!4}
= 3\!\left(\frac{{G^{*}}^2}{4\varpi}\right)^{\!4}
= \frac{3\,{G^{*}}^8}{256\,\varpi^4}\,.\qedhere
\]
\end{proof}

\begin{remark}\label{rem:2d-ontic}
The full two-dimensional lattice structure at the Gaussian point is
controlled entirely by $\varpi$ and~${G^{*}}$.  Our one-dimensional sum~$S$
is a restriction of this structure, and the restriction to a single row is
what introduces the circular constant~$\pi$ as an intermediary through the
bilateral-sum formula~\eqref{eq:bilateral}.
\end{remark}

%=======================================================================
\section{The Epstein--Lemniscatic Connection}\label{sec:epstein}
%=======================================================================

\subsection{The Epstein Factorisation in Ontic Variables}\label{sec:epstein-ontic}

We now close the loop between the Epstein zeta function and the lemniscatic
constants.  The classical factorisation of $Z_{\Z[i]}(s)$ through
Dirichlet series gives:

\begin{theorem}[Epstein factorisation]\label{thm:epstein-factor}
\begin{equation}\label{eq:epstein-factor}
Z_{\Z[i]}(s) \;=\; 4\,\zeta(s)\,L(s,\chi_{-4})\,,
\end{equation}
where $\zeta(s)$ is the Riemann zeta function and $L(s,\chi_{-4})$ is the
Dirichlet $L$-function attached to the non-principal character modulo~$4$
(the Kronecker symbol $\chi_{-4}(n)=\bigl(\frac{-4}{n}\bigr)$).
\end{theorem}
\begin{proof}
The number of representations $r_2(k)=\#\{(m,n)\in\Z^2:m^2+n^2=k\}$
satisfies $r_2(k)=4\sum_{d\mid k}\chi_{-4}(d)$, so
\[
Z_{\Z[i]}(s)=\sum_{k=1}^{\infty}r_2(k)\,k^{-s}
=4\sum_{k=1}^{\infty}\biggl(\sum_{d\mid k}\chi_{-4}(d)\biggr)k^{-s}
=4\,\zeta(s)\,L(s,\chi_{-4})\,.\qedhere
\]
\end{proof}

At $s=2$, each factor acquires a specific identity:

\begin{corollary}[Ontic form of $Z_{\Z[i]}(2)$]\label{cor:Z2-ontic}
\begin{equation}\label{eq:Z2-ontic}
Z_{\Z[i]}(2) \;=\; 4\,\zeta(2)\,G
\;=\;\frac{32\,\varpi^4}{3\,{G^{*}}^4}\;\cdot\; G\,,
\end{equation}
where $G=L(2,\chi_{-4})\approx 0.915966$ is Catalan's constant.
\end{corollary}
\begin{proof}
Substitute $\zeta(2)=\pi^2/6$ and $\pi=4\varpi^2/{G^{*}}^2$:
\[
\zeta(2)=\frac{\pi^2}{6}=\frac{16\varpi^4}{6\,{G^{*}}^4}
=\frac{8\varpi^4}{3\,{G^{*}}^4}\,.
\]
Then $Z_{\Z[i]}(2)=4\cdot\dfrac{8\varpi^4}{3\,{G^{*}}^4}\cdot G
=\dfrac{32\,\varpi^4\,G}{3\,{G^{*}}^4}$.
\end{proof}

\begin{remark}[The irreducible arithmetic residue]\label{rem:catalan}
Catalan's constant
\[
G \;=\; \sum_{n=0}^{\infty}\frac{(-1)^n}{(2n+1)^2}
\;=\; 1-\frac{1}{9}+\frac{1}{25}-\frac{1}{49}+\cdots
\;\approx\;0.915966
\]
encodes the distribution of Gaussian primes: it measures the excess of
primes $p\equiv 1\pmod{4}$ (which split in~$\Z[i]$) over primes
$p\equiv 3\pmod{4}$ (which remain inert), weighted by $p^{-2}$.  This
arithmetic content is intrinsic to the \emph{two-dimensional} Gaussian
lattice and cannot be expressed in terms of $\varpi$ and~${G^{*}}$ alone.
It is the irreducible arithmetic residue that distinguishes the norm
projection $|z|^{-2s}$ from the holomorphic projection $z^{-2s}$.
\end{remark}

%-----------------------------------------------------------------------
\subsection{The Row Decomposition}\label{sec:row-decomp}
%-----------------------------------------------------------------------

The Epstein zeta function admits a natural decomposition by rows
$n=\text{const}$:
\begin{equation}\label{eq:row-decomp}
Z_{\Z[i]}(2) \;=\; 2\,\zeta(4)
+ 2\sum_{n=1}^{\infty}\sum_{m=-\infty}^{\infty}\frac{1}{(m^2+n^2)^2}\,,
\end{equation}
where $2\,\zeta(4)=\sum'_m m^{-4}$ is the $n=0$ row and the factor of~$2$
in front of the sum accounts for the $n$ and $-n$ rows.

Each bilateral row sum can be evaluated by differentiating~\eqref{eq:bilateral}
with respect to~$a$ and setting $a=n^2$:
\begin{equation}\label{eq:bilateral-row}
\sum_{m=-\infty}^{\infty}\frac{1}{(m^2+n^2)^2}
\;=\;\frac{\pi^2\csch^2(n\pi)}{2n^2}+\frac{\pi\coth(n\pi)}{2n^3}\,.
\end{equation}
In ontic variables, each row depends on $n$ through $\coth(4n\ell)$
and $\csch(4n\ell)$.

\begin{proposition}[Row decomposition in ontic form]\label{prop:row-ontic}
\begin{equation}\label{eq:row-ontic}
Z_{\Z[i]}(2) = \frac{2^8\,\varpi^8}{45\,{G^{*}}^8}
+ 2\sum_{n=1}^{\infty}\left[
\frac{8\varpi^4\csch^2(4n\ell)}{n^2\,{G^{*}}^4}
+ \frac{2\varpi^2\coth(4n\ell)}{n^3\,{G^{*}}^2}\right],
\end{equation}
where $\ell=\varpi^2/{G^{*}}^2$ and the leading term uses
$\zeta(4)=\pi^4/90=2^8\varpi^8/(90\,{G^{*}}^8)$.
\end{proposition}

\begin{remark}[Why the row restriction eliminates $G$]\label{rem:row-kills-G}
Each individual row sum~\eqref{eq:bilateral-row} is expressible purely in
terms of $\varpi$ and~${G^{*}}$ (via $\pi=4\varpi^2/{G^{*}}^2$).  Catalan's
constant enters only when we sum the row contributions and equate them to
the Dirichlet-series factorisation $4\,\zeta(2)\,L(2,\chi_{-4})$.  The row
restriction to $n=1$ that yields our sum~$S$ bypasses this global
arithmetic content entirely: it uses only the Poisson summation formula
over~$\Z$, whose kernel is the Fourier transform of $1/(x^2+1)^2$ --- a
smooth function controlled by the single period~$\varpi$.

This explains why $S$ admits a purely ontic closed form
(Theorem~\ref{thm:ftd}), while the full $Z_{\Z[i]}(2)$ does not: the
two-dimensional lattice sum encodes arithmetic information (the
distribution of sums of two squares) that the one-dimensional restriction
projects away.
\end{remark}

%-----------------------------------------------------------------------
\subsection{Two Projections of One Lattice}\label{sec:two-projections}
%-----------------------------------------------------------------------

The Gaussian lattice $\Z[i]$ supports two natural ``projections'' at
weight~4, each revealing different content:

\medskip
\begin{center}
\begin{tabular}{@{}llll@{}}
\toprule
\textbf{Projection} & \textbf{Object} & \textbf{FTD form} & \textbf{Ontic purity}\\
\midrule
Holomorphic ($z^{-4}$) & $G_4(i),\;E_4(i)$
  & $3\,{G^{*}}^8/(256\,\varpi^4)$
  & Fully ontic\\[3pt]
Norm ($|z|^{-4}$) & $Z_{\Z[i]}(2)$
  & $32\,\varpi^4\,G/(3\,{G^{*}}^4)$
  & Ontic $\times$ arithmetic\\[3pt]
Row restriction & $S$
  & $\dfrac{2+4\ell\coth 4\ell+16\ell^2\csch^2 4\ell}{4}$
  & Fully ontic\\
\bottomrule
\end{tabular}
\end{center}
\medskip

The holomorphic projection $z^{-4}$ respects the four-fold symmetry of
$\Z[i]$ and yields a holomorphic modular form whose value at $\tau=i$ is
determined by the lattice periods $\varpi$ and $i\varpi$ alone.  The norm
projection $|z|^{-4}$ discards phase information and introduces the
arithmetic of $\Z[i]$ through the representation function $r_2(k)$, which
carries Catalan's constant as its weight-2 residue.

\begin{remark}\label{rem:bridge-is-Gstar}
In both projections, the lemniscatic constants $\varpi$ and ${G^{*}}$ control
the geometric structure of the lattice.  The \emph{arithmetic} factor $G$
appears only when we ask a number-theoretic question (``how many lattice
points lie on the circle $|z|^2=k$?'') rather than a geometric one
(``what are the periods?'').  The bridge constant ${G^{*}}$ mediates the
geometric content; Catalan's constant encodes the arithmetic residue.
\end{remark}

\subsection{Numerical Verification}\label{sec:numerical-ext}

\begin{center}
\begin{tabular}{lll}
\toprule
\textbf{Quantity} & \textbf{Value} & \textbf{Method}\\
\midrule
$S$ (direct summation, $10^5$ terms) & 1.306836975422908 & Numerical\\
$S$ (closed form, Eq.\,\ref{eq:result}) & 1.306836975422909 & Analytic\\
$S$ (FTD form, Eq.\,\ref{eq:ftd-form}) & 1.306836975422909 & Analytic\\
$\pi$ via $4\varpi^2/{G^{*}}^2$ & 3.141592653589793 & Verified\\
$E_4(i)$ ($q$-expansion, 500 terms) & 1.455762892268709 & Numerical\\
$3{G^{*}}^8/(256\varpi^4)$ & 1.455762892268710 & Analytic\\
\midrule
$\zeta(2)$ & 1.644934066848226 & Standard\\
$8\varpi^4/(3{G^{*}}^4)$ & 1.644934066848226 & Ontic\\
Catalan $G = L(2,\chi_{-4})$ & 0.915965594177219 & Standard\\
$Z_{\Z[i]}(2)$ via $4\,\zeta(2)\,G$ & 6.026812039691940 & Factored\\
$32\varpi^4 G/(3{G^{*}}^4)$ & 6.026812039691940 & Ontic\\
$Z_{\Z[i]}(2)$ via row decomposition & 6.026812039691940 & Row sum\\
$n\!=\!1$ bilateral row & $1.613673950845818$ & Analytic\\
$2S-1$ & $1.613673950845818$ & Verified\\
\bottomrule
\end{tabular}
\end{center}

%=======================================================================
\section{Structure of the Discrete--Continuous Bridge}\label{sec:bridge-structure}
%=======================================================================

The derivation passes through three distinct bridges, each moving from a
more discrete to a more continuous description:

\medskip
\begin{center}
\begin{tabular}{clll}
\toprule
\textbf{Step} & \textbf{From} & \textbf{To} & \textbf{Mechanism}\\
\midrule
1 & Integer lattice $\Z$ & Hyperbolic function of $a$ & Bilateral sum / Poisson\\
2 & Fixed sum at $a=1$ & Continuous family in $a$ & Parameterisation\\
3 & First-order sum & Squared denominator & Differentiation $-d/da$\\
\bottomrule
\end{tabular}
\end{center}
\medskip

The final answer decomposes into three additive contributions (within the
numerator):

\medskip
\begin{center}
\begin{tabular}{llcl}
\toprule
\textbf{Term} & \textbf{Expression} & \textbf{Share} & \textbf{Interpretation}\\
\midrule
Constant    & $2/4=0.500$                   & 38.3\% & Pure discrete residue\\
Linear in $\ell$ & $4\ell\coth(4\ell)/4$  & 60.3\% & Bridge term (${G^{*}}$ mediation)\\
Quadratic in $\ell$ & $16\ell^2\csch^2(4\ell)/4$ & 1.4\% & Curvature correction\\
\bottomrule
\end{tabular}
\end{center}
\medskip

The ``constant'' term $2/4=1/2$ is the part that would survive if the
hyperbolic functions were replaced by their large-argument limits
($\coth\to 1$, $\csch\to 0$), corresponding to $\ell\to\infty$ or
equivalently $q\to 0$.  It represents the purely discrete contribution:
the value the sum would approach if the lattice spacing were sent to zero
relative to the period.

The ``linear'' term $4\ell\coth(4\ell)/4$ carries 60.3\% of the
answer and is dominated by the ratio $\varpi^2/{G^{*}}^2$.  This is the
bridge term proper: it encodes the interaction between the discrete lattice
and the continuous exponential.

The ``quadratic'' correction $16\ell^2\csch^2(4\ell)/4$ contributes
only 1.4\%.  Its suppression is a direct consequence of the smallness of
the nome $q=e^{-8\ell}\approx 1.87\times 10^{-3}$, whose magnitude is
set by $\ell=\varpi^2/{G^{*}}^2$ alone.

%=======================================================================
\section{The Genus-0 Projection}\label{sec:genus}
%=======================================================================

The bilateral sum formula~\eqref{eq:bilateral} ultimately rests on the
partial-fraction expansion of $\pi\cot(\pi z)$, which is a genus-0
(rational curve) identity.  The hyperbolic functions $\coth$ and $\csch$
are all genus-0 objects: they parametrise the Riemann sphere $\mathbb{P}^1$.

The Gaussian lattice $\Z[i]$, however, is a genus-1 object.  Its natural
functions are the Weierstrass $\wp$-function and its derivatives, whose
periods are $\varpi$ and~$i\varpi$, and whose invariants are expressible purely
in terms of $\varpi$ and~${G^{*}}$.  (For the square lattice, $g_3=0$ by
four-fold symmetry, and $g_2$ is determined by~$\varpi$ alone.)

The fact that the answer can be written in terms of $\coth$ and $\csch^2$
of~$\pi$ is a consequence of the one-dimensional restriction collapsing the
elliptic (genus-1) structure down to the circular (genus-0) one.
Restricting to a single row $n=1$ in the Epstein sum replaces the
doubly-periodic lattice with a singly-periodic one, and single periodicity
is the domain of trigonometric and hyperbolic functions.

In the FTD reading:
\begin{quote}
$\pi$ appears in this formula not because the sum is ``about'' circles,
but because we chose to evaluate a lemniscatic object using circular tools.
The ontic content is $\ell=\varpi^2/{G^{*}}^2$.  Everything else is
projection.
\end{quote}

%=======================================================================
\section{The Dimensional Ladder}\label{sec:ladder}
%=======================================================================

The genus-0 projection explains why $\pi$ appears in the answer: we
evaluated a two-dimensional (genus-1) object using one-dimensional
(genus-0) tools.  But this raises a deeper question: why $\pi$ at all?

Euler's 1735 derivation of the Basel problem provides the answer.  His
method was to observe that $\sin(\pi z)/(\pi z)$ vanishes at every
nonzero integer, and to treat it---by analogy with finite
polynomials---as an infinite product over its roots:
\begin{equation}\label{eq:euler-product}
\frac{\sin(\pi z)}{\pi z}
\;=\; \prod_{n=1}^{\infty}\!\left(1 - \frac{z^2}{n^2}\right).
\end{equation}
Expanding the right side and comparing the $z^2$ coefficient with the
Taylor series $\sin(\pi z)/(\pi z) = 1 - \pi^2 z^2/6 + \cdots$ gives
$\sum 1/n^2 = \pi^2/6$.

Now substitute $\pi = 4\ell$, where $\ell = \varpi^2/{G^{*}}^2$ is the
ontic ratio:
\[
\frac{\sin(4\ell z)}{4\ell z}
\;=\; \prod_{n=1}^{\infty}\!\left(1 - \frac{z^2}{n^2}\right),
\qquad
\sum_{n=1}^{\infty}\frac{1}{n^2}
\;=\; \frac{(4\ell)^2}{6}
\;=\; \frac{8\ell^2}{3}
\;=\; \frac{8\varpi^4}{3{G^{*}}^4}\,.
\]
The right side of the product formula is unchanged---it depends only on
the integers.  The left side has replaced $\pi$ with $4\ell$.  The
substitution is algebraically trivial, but it reveals the structural
origin: the number $\pi^2/6$ is the zeta function of the
one-dimensional integer lattice~$\Z$ evaluated at $s = 2$.  The zeros
of $\sin(\pi z)$ are the lattice points of~$\Z$.  The Basel sum is a
\emph{lattice sum}.

This observation generalises.  Each dimension has its own lattice, its
own fundamental period, and its own zeta function:

\medskip
\begin{center}
\begin{tabular}{cllll}
\toprule
$D$ & Lattice & Period & Lattice zeta at $s=2$ & Product formula\\
\midrule
1 & $\Z$ & $\pi = 4\ell$
  & $\displaystyle\zeta(2) = \frac{\pi^2}{6} = \frac{8\varpi^4}{3{G^{*}}^4}$
  & $\displaystyle\frac{\sin\pi z}{\pi z} = \prod\!\left(1-\frac{z^2}{n^2}\right)$\\[10pt]
2 & $\Z[i]$ & $\varpi$
  & $\displaystyle Z_{\Z[i]}(2) = 4\,\zeta(2)\,L(2,\chi_{-4})$
  & Weierstrass $\sigma(z;\Z[i])$\\[10pt]
3 & $\Z^3$ & ${G^{*}}$
  & $\displaystyle W_3 = \frac{{G^{*}}^2}{2\pi} = \frac{\Gamma(1/4)^4}{4\pi^3}$
  & Watson integral\\
\bottomrule
\end{tabular}
\end{center}
\medskip

In $D = 1$, the ``Euler move''---extending finite polynomial
factorisations to infinite products---discovers $\pi$.  The same move
applied to the Gaussian lattice $\Z[i]$ in $D = 2$ discovers $\varpi$:
the Weierstrass $\sigma$-function of $\Z[i]$ is the doubly-periodic
analogue of $\sin(\pi z)$, and its periods are $\varpi$ and $i\varpi$.  In
$D = 3$, the lattice Green's function of $\Z^3$---the Watson
integral---discovers ${G^{*}}$, which absorbs the geometric scaling
between the lemniscatic and circular domains.

Each lattice period is a ``circle constant'' for its dimension:
\begin{itemize}
\item $\pi$ measures how far you walk around a circle (the boundary of
  the unit ball in $\R^1$ cross-section).
\item $\varpi$ measures how far you walk around a lemniscate (the
  fundamental domain boundary of $\Z[i]$).
\item ${G^{*}}$ measures the lattice Green's function of $\Z^3$---the
  propagator of the cubic lattice itself.
\end{itemize}

The reason $\pi$ appears ubiquitously in mathematics is not that circles
are fundamental.  It is that $\Z$---the one-dimensional integer
lattice---is the simplest lattice, and most classical formulas are
computed over it.  When we compute over $\Z[i]$, we find $\varpi$.  When
we compute over $\Z^3$, we find ${G^{*}}$.

The connection to FTD's dimension selection is immediate.  The master
quadratic $x^2 - 16{G^{*}}_D^2\,x + 16{G^{*}}_D^3 = 0$, applied to the
Watson integral $W_D$ in each dimension, gives $\lfloor x_-^{(D)}\rfloor = D$
only for $D = 3$~\cite{lifecycle}.  Among the infinite family of lattice
periods $(\pi,\; \varpi,\; {G^{*}},\; {G^{*}}_4,\; \ldots)$, only the
$D = 3$ member produces a self-consistent framework in which the number
of color charges equals the spatial dimension.

The dimensional ladder is summarized in the following table:

\begin{center}
\begin{tabular}{cclll}
\toprule
$D$ & Lattice & Period & Zeta function & $\lfloor x_-^{(D)}\rfloor \stackrel{?}{=} D$ \\
\midrule
1 & $\Z$ & $\pi$ & $\zeta(2) = \pi^2/6$ & N/A \\
2 & $\Z[i]$ & $\varpi$ & $G_4(i) \propto \varpi^4$ & \textbf{No} ($8 \neq 2$) \\
3 & $\Z^3$ & ${G^{*}}$ & $W_3 = {{G^{*}}}^2/(2\pi)$ & \textbf{Yes} ($3 = 3$) \\
\bottomrule
\end{tabular}
\end{center}

\noindent Each dimension has its own ``circle constant.'' $\pi$ is the period of $\Z$. $\varpi$ is the period of $\Z[i]$. ${G^{*}}$ is the period of $\Z^3$. The relation $\pi = 4\varpi^2/{{G^{*}}}^2$ expresses the lowest-dimensional constant through the higher ones---$\pi$ is the projection of lemniscatic geometry onto a line.

Euler's Basel problem is the one-dimensional shadow of a three-dimensional structure.

\begin{proposition}[Uniqueness of $D = 3$]
Among dimensions $D = 1, 2, \ldots, 6$, the gap equation $x^2 - 16G_D^{*2}x + 16G_D^{*3} = 0$ (where $G_D^* = \sqrt{2\pi W_D}$ and $W_D$ is the Watson integral for $\mathbb{Z}^D$) satisfies $\lfloor x_-(D) \rfloor = D$ if and only if $D = 3$.
\end{proposition}

\begin{remark}
The monotonicity of $x_-(D)$ for $D \geq 4$ ensures this uniqueness extends to all $D \geq 1$.
\end{remark}

%=======================================================================
\section{Summary of Rewriting}\label{sec:summary}
%=======================================================================

\begin{center}
\begin{tabular}{lll}
\toprule
\textbf{Quantity} & \textbf{Standard Form} & \textbf{FTD (Ontic) Form}\\
\midrule
Argument of $\coth/\csch$ & $\pi$ & $4\varpi^2/{G^{*}}^2$\\
Exponential scale & $e^{\pi}$ & $e^{4\varpi^2/{G^{*}}^2}$\\
Lattice nome & $e^{-2\pi}$ & $e^{-8\varpi^2/{G^{*}}^2}$\\
$E_4(i)$ & $3(\varpi/\pi)^4$ & $3{G^{*}}^8/(256\varpi^4)$\\
Prefactor $\pi^2$ & $\pi^2$ & $16\varpi^4/{G^{*}}^4$\\
$\zeta(2)$ & $\pi^2/6$ & $8\varpi^4/(3{G^{*}}^4)$\\
Euler product & $\sin(\pi z)/(\pi z) = \prod(1-z^2/n^2)$ & $\sin(4\ell z)/(4\ell z) = \prod(1-z^2/n^2)$\\
$Z_{\Z[i]}(2)$ & $4\,\zeta(2)\,G$ & $32\varpi^4 G/(3{G^{*}}^4)$\\
Catalan $G$ & $L(2,\chi_{-4})$ & (irreducible arithmetic constant)\\
\bottomrule
\end{tabular}
\end{center}
\medskip

\noindent\textbf{Conclusion.}
The discrete--continuous bridge in the evaluation of
$\sum 1/(n^2+1)^2$ is real, but the conventional language in terms of
$\pi$ obscures its structural origin.  The bridge is not~$\pi$.  The
bridge is~${G^{*}}$.

The Epstein zeta factorisation~\eqref{eq:epstein-factor} reveals an
additional layer: the full two-dimensional lattice sum carries an
irreducible arithmetic factor (Catalan's constant~$G$) that encodes the
distribution of Gaussian primes.  This factor is eliminated by the row
restriction that produces our one-dimensional sum~$S$, explaining why~$S$
admits a purely ontic closed form while the full $Z_{\Z[i]}(2)$ does not.

%=======================================================================
\section*{Epistemic Classification}
%=======================================================================

Following FTD conventions:

\begin{itemize}
\item \textbf{THEOREM} (exact algebraic): Equations~\eqref{eq:result},
\eqref{eq:ftd-form}, \eqref{eq:E4}, \eqref{eq:epstein-factor},
\eqref{eq:Z2-ontic}, and Lemma~\ref{lem:pi}.
These are mathematically exact identities, verifiable by direct computation.

\item \textbf{THEOREM} (standard number theory):
The factorisation $Z_{\Z[i]}(2)=4\,\zeta(2)\,L(2,\chi_{-4})$ and the
identification of Catalan's constant as the irreducible arithmetic
residue.  The row-restriction elimination of the arithmetic factor follows
from Poisson summation.

\item \textbf{SELECTION PRINCIPLE} (pending proof):
The claim that $\varpi$ and ${G^{*}}$ are ontologically prior to~$\pi$.
This is a well-motivated structural claim supported by the lemniscatic
lattice framework but not provable within conventional mathematics, which
treats all equivalent parametrisations as interchangeable.

\item \textbf{HYPOTHESIS} (empirical fit):
The identification of ${G^{*}}$ as the mediator of all
discrete--continuous transitions in physics.  This is a broader FTD claim
that this example illustrates but does not prove.
\end{itemize}

%=======================================================================
\begin{thebibliography}{9}

\bibitem{zucker77}
I.\,J.~Zucker,
``The evaluation in terms of $\Gamma$-functions of the periods of elliptic
curves admitting complex multiplication,''
\textit{Math.\ Proc.\ Camb.\ Phil.\ Soc.} \textbf{82} (1977), 111--118.

\end{thebibliography}

\end{document}
